% config/macros.tex — macros de trazabilidad y recuadros del libro

% --- Trazabilidad contra el programa oficial (§4) ---
% Uso: \progtag{1.1.1} coloca una etiqueta \label{prog:1.1.1} invisible en el margen,
% usada por el script de verificación de cobertura (docs/cobertura.md).
\newcommand{\progtag}[1]{\label{prog:#1}}

% --- Marcadores de pendientes (§5.2, §9) — visibles en el PDF y listados aparte ---
\newcommand{\pendiente}[1]{%
  \textcolor{colorAtencion}{\textbf{[PENDIENTE: #1]}}%
}
\newcommand{\datoPendiente}[1]{%
  \textcolor{colorAtencion}{\textbf{[DATO PENDIENTE DE VERIFICAR: #1]}}%
}

% --- Recuadros tcolorbox (§11) ---
\usepackage[most]{tcolorbox}
\tcbuselibrary{skins,breakable}

\newtcolorbox{cajaDefinicion}[1][]{%
  colback=colorDefinicion!5, colframe=colorDefinicion, breakable,
  fonttitle=\bfseries, title={Definición}, #1}

\newtcolorbox{cajaEjemplo}[1][]{%
  colback=colorEjemplo!5, colframe=colorEjemplo, breakable,
  fonttitle=\bfseries, title={Ejemplo}, #1}

\newtcolorbox{cajaAtencion}[1][]{%
  colback=colorAtencion!5, colframe=colorAtencion, breakable,
  fonttitle=\bfseries, title={Atención}, #1}

\newtcolorbox{cajaFrontera}[1][]{%
  colback=colorFrontera!5, colframe=colorFrontera, breakable,
  fonttitle=\bfseries, title={Frontera}, #1}

\newtcolorbox{cajaEtica}[1][]{%
  colback=colorEtica!5, colframe=colorEtica, breakable,
  fonttitle=\bfseries, title={Ética y responsabilidad con los datos}, #1}

\newtcolorbox{cajaLaboratorio}[1][]{%
  colback=colorLaboratorio!5, colframe=colorLaboratorio, breakable,
  fonttitle=\bfseries, title={Laboratorio MATLAB}, #1}

% Ficha del capítulo (§8.1): recuadro neutro, administrativo — no compite
% visualmente con los recuadros de contenido (Definición, Frontera, Ética...).
\newtcolorbox{cajaFicha}[1][]{%
  colback=grisIPN!6, colframe=grisIPN, breakable, sharp corners,
  boxrule=0.6pt, left=8pt, right=8pt, top=4pt, bottom=4pt,
  fonttitle=\bfseries, title={Ficha del capítulo}, #1}

% --- Vigencia de temas selectos (§5.7) ---
\newcommand{\vigente}[1]{\textit{(Vigente a #1)}}

% --- Caso ilustrativo ficticio (§5.6) ---
\newenvironment{casoFicticio}{%
  \begin{cajaEjemplo}[title={Caso ilustrativo (ficticio)}]%
}{%
  \end{cajaEjemplo}%
}
